\documentclass[11pt]{article}
\usepackage[final]{acl}
\usepackage{times}
\usepackage{latexsym}
\usepackage{fontspec}
\usepackage{polyglossia}
\setdefaultlanguage{russian}
\setotherlanguage{english}
\setmainfont{Times New Roman}
\newfontfamily\cyrillicfont{Times New Roman}
\usepackage{microtype}
\usepackage{inconsolata}
\usepackage{graphicx}

\title{Ранжирование BM25 и оценка качества поиска}
\author{Игорь Пивнев, ИТМО}

\begin{document}
\maketitle

\begin{abstract}
В работе исследуются классические методы информационного поиска на коллекции
WikIR (подмножество en1k). Рассматриваются модели TF-IDF и BM25, а также влияние
морфологической обработки (стемминг и лемматизация) на качество поиска. Проводится
оценка с использованием стандартных метрик (Precision@k, MAP, nDCG), а также
анализ результатов и подбор параметров BM25.
\end{abstract}

\section{Постановка задачи}

Целью работы является исследование качества поиска документов по текстовым
запросам с использованием классических моделей ранжирования. 

Решаются следующие задачи:
\begin{itemize}
    \item анализ набора запросов и релевантных документов;
    \item реализация поиска с использованием TF-IDF и BM25;
    \item сравнение различных вариантов предобработки текста;
    \item оценка качества поиска с использованием стандартных IR-метрик;
    \item анализ полученных результатов;
    \item подбор параметров модели BM25.
\end{itemize}

\section{Описание данных}

В работе используется коллекция WikIR (подмножество en1k), содержащая:
\begin{itemize}
    \item документы (тексты статей Wikipedia);
    \item поисковые запросы;
    \item оценки релевантности (qrels).
\end{itemize}

Основные характеристики тестового набора данных:
\begin{itemize}
    \item количество запросов: \textbf{100};
    \item средняя длина запроса: \textbf{2.07} слов;
    \item среднее число релевантных документов на запрос: \textbf{44.35}.
\end{itemize}

Запросы в коллекции получены автоматически на основе структуры Wikipedia, а
релевантность определяется по связанным документам, что делает задачу приближенной
к реальным условиям информационного поиска. Для ускорения экспериментов из полной
коллекции была взята случайная выборка из 50\,000 документов.

\section{Методы и инструменты}

В работе использовались две основные модели ранжирования:

\textbf{TF-IDF} — векторная модель, в которой важность терма определяется его
частотой в документе и обратной частотой в коллекции. Для вычислений применялась
реализация из \texttt{sklearn.feature\_extraction.text.TfidfVectorizer}.

\textbf{BM25} — вероятностная модель ранжирования, учитывающая насыщение частоты
терма и длину документа:
\begin{equation}
\text{score}(q, d) = \sum_{t \in q} IDF(t) \cdot \frac{f(t,d)(k_1+1)}{f(t,d)+k_1(1-b+b\cdot \frac{|d|}{avgdl})}
\end{equation}
Использовалась реализация из библиотеки \texttt{rank\_bm25}.

Рассматривались три варианта обработки текста:
\begin{itemize}
    \item оригинальный текст (токенизация по пробелам);
    \item стемминг (Porter Stemmer из NLTK);
    \item лемматизация (модель \texttt{en\_core\_web\_sm} библиотеки spaCy).
\end{itemize}

Для оценки качества использовались метрики:
\begin{itemize}
    \item Precision@1, Precision@10, Precision@20;
    \item MAP (Mean Average Precision);
    \item nDCG@20 (с учетом градуированной релевантности).
\end{itemize}

Оценка проводилась с использованием библиотеки \texttt{ir-measures}. Также
измерялось время выполнения запросов для всех моделей.

\section{Результаты и интерпретация}

Результаты экспериментов представлены в таблице~\ref{tab:results}.

\begin{table}[h]
\centering
\begin{tabular}{lccccc}
\hline
Модель & P@1 & P@10 & MAP & nDCG@20 \\
\hline
TF-IDF & 0.23 & 0.062 & 0.0344 & 0.0950 \\
BM25 & 0.27 & 0.069 & 0.0429 & 0.1107 \\
BM25 + stem & 0.27 & 0.067 & 0.0410 & 0.1108 \\
BM25 + lemma & 0.27 & 0.067 & 0.0417 & 0.1102 \\
\hline
\end{tabular}
\caption{Сравнение моделей}
\label{tab:results}
\end{table}

Время выполнения запросов (секунд):
\begin{itemize}
    \item TF-IDF: \textbf{1.86};
    \item BM25: \textbf{4.49};
    \item BM25 со стеммингом: \textbf{4.36};
    \item BM25 с лемматизацией: \textbf{4.36}.
\end{itemize}

Основные наблюдения:
\begin{itemize}
    \item BM25, как правило, показывает более высокое качество по сравнению с
    TF-IDF за счет учета насыщения частоты и нормализации длины документа.
    \item Морфологическая обработка может как улучшать, так и ухудшать результаты.
    В данном эксперименте лемматизация дала небольшое улучшение по большинству
    метрик, тогда как стемминг привел к незначительному падению качества.
    \item Лемматизация требует больше вычислительных ресурсов, но обеспечивает
    более точное приведение слов к нормальной форме, что полезно для языков с
    развитой морфологией (в нашем случае английский).
\end{itemize}

\section{Анализ результатов}

\subsection{Сложные и простые запросы}
Были выявлены запросы с низким и высоким значением MAP. Простые запросы, как
правило, содержат редкие и специфичные термы (например, имена собственные или узкие
понятия), тогда как сложные — короткие и неоднозначные (например, общие слова типа
"history", "science"). Примеры самых сложных и самых простых запросов приведены в коде.

\subsection{Влияние длины запроса}
Более длинные запросы обычно дают лучшее качество, так как содержат больше
информации о намерении пользователя. Корреляция между длиной запроса и AP составила
\textbf{0.1278}. Это ожидаемый результат, поскольку длинные запросы легче
сопоставить с релевантными документами.

\subsection{Влияние морфологической обработки}
Стемминг иногда приводит к потере смысла (слишком агрессивное обрезание, например,
объединение слов с разным значением). Лемматизация сохраняет семантику, но требует
больше вычислений. В нашем эксперименте лемматизация показала себя более стабильно,
чем стемминг.

\subsection{Сравнение метрик}
Разные метрики дают различную картину качества:
\begin{itemize}
    \item Precision@k чувствительна к первым позициям;
    \item MAP учитывает весь ранжированный список;
    \item nDCG учитывает градуированную релевантность.
\end{itemize}
Использование одной метрики недостаточно для полной оценки качества. Например,
модель может иметь высокий P@10, но низкий MAP, если релевантные документы
распределены по списку неравномерно.

\subsection{Ошибки модели}
Среди топовых документов встречаются нерелевантные из-за:
\begin{itemize}
    \item совпадения по ключевым словам без учета контекста;
    \item синонимии и полисемии;
    \item ограниченности bag-of-words представления.
\end{itemize}
Примеры ошибок: документ, содержащий слово "bank" в значении "берег реки", выходит
на запрос "bank" в финансовом смысле.

\subsection{Подбор параметров BM25}
Был выполнен перебор параметров $k_1$ и $b$ на валидационной выборке. Лучшие значения:
\begin{itemize}
    \item $k_1 = $ \textbf{0.5};
    \item $b = $ \textbf{0.5}.
\end{itemize}

Подбор параметров позволил \textbf{незначительно изменить} качество модели, что
подтверждает важность настройки гиперпараметров.

\section*{Ограничения и возможные улучшения}

Основные ограничения работы:
\begin{itemize}
    \item использование простых bag-of-words моделей, не учитывающих порядок слов
    и семантику;
    \item отсутствие семантического понимания текста;
    \item зависимость от качества предобработки (стемминг/лемматизация);
    \item выборка из 50\,000 документов, что может не полностью отражать свойства
    всей коллекции.
\end{itemize}

Возможные улучшения:
\begin{itemize}
    \item использование нейронных моделей (BERT, ColBERT) для учета контекста;
    \item query expansion (расширение запроса синонимами или связанными терминами);
    \item обучение ранжирующих моделей (learning-to-rank) на основе признаков;
    \item использование семантических эмбеддингов (например, SBERT) для поиска по смыслу;
    \item применение более сложных методов предобработки, включая разрешение омонимии.
\end{itemize}

\end{document}
